两片无限理想导体平板位于 $x=0$ 与 $x=a$，$a>0$，板间填充介电常数 $\epsilon>0$、磁导率 $\mu>0$ 的无损均匀介质。波导沿 $+z$ 无限延伸，研究 $y$ 偏振的最低阶 TE 型模式，其电场按复表示为
\[
E_y(x,z,t)=E_0\sin(\pi x/a)e^{i(k_z z-\omega t)}
\]
且其他电场分量为零，$\omega>0$。取 $e^{-i\omega t}$ 时间相位约定。
(1) 从介质中的 Helmholtz 方程和理想导体切向电场边界条件推导 $k_z(\omega)$ 及截止频率；(2) 求高于截止时的相速度、群速度与 TE 波阻抗，解释相速度大于介质光速为何不违反因果律；(3) 从 Maxwell–Faraday 公式求磁场的 $H_x,H_z$ 分量，并计算横截面平均坡印廷功率流；(4) 对低于截止的频率求沿波导方向的倏逝长度，讨论此时为何没有向无限远输运的平均功率。
